% =====================================================================
% P8 — The Forced Internationalization Penalty (FIP): a monotone-negative
%       I-P boundary condition in Pacific Small Island Developing States
% WORKING DRAFT submission framework (Springer-compatible).
%
% COMPILE:
%   pdflatex p8_sids_fip.tex  (twice) ; bibtex p8_sids_fip ; pdflatex (x2)
% NOTE: No pdflatex was available in the build environment. Uses standard
%   `article` class (Springer sn-jnl.cls is not bundled here). Swap the
%   class/title block for sn-jnl when submitting to a Springer journal.
%
% INTEGRITY NOTE: Every numeric value is verbatim from the dissertation
%   (chuong_4_ket_qua_vi.md sec 4.7 / 4.9; methods from phu_luc_A). Full
%   regression tables absent from the dissertation are explicit
%   PLACEHOLDERS. No new coefficients invented.
% =====================================================================

\documentclass[12pt,a4paper]{article}
\usepackage[margin=2.5cm]{geometry}
\usepackage{setspace}\onehalfspacing
\usepackage{iftex}
\ifPDFTeX
  \usepackage[utf8]{inputenc}\usepackage[T1]{fontenc}\usepackage{lmodern}
\else
  \usepackage{fontspec}
\fi
\usepackage{amsmath,amssymb,amsfonts,mathtools}
\usepackage{booktabs}\usepackage{array}\usepackage{tabularx}
\usepackage{graphicx}\usepackage{float}
\usepackage{caption}\captionsetup{font=small,labelfont=bf}
\usepackage{microtype}\usepackage{textcomp}
\usepackage[hidelinks,breaklinks=true]{hyperref}
\usepackage[round]{natbib}\bibliographystyle{apalike}

\usepackage{filecontents}
\begin{filecontents*}{p8refs.bib}
@article{deaton1985,
  author={Deaton, Angus}, title={Panel data from a time series of cross-sections},
  journal={Journal of Econometrics}, volume={30}, number={1--2}, pages={109--126}, year={1985}}
@incollection{verbeek2008,
  author={Verbeek, Marno}, title={Pseudo-panels and repeated cross-sections},
  booktitle={The Econometrics of Panel Data}, edition={3rd}, pages={369--383},
  publisher={Springer}, year={2008}}
@article{lindmehlum2010,
  author={Lind, Jo Thori and Mehlum, Halvor},
  title={With or without {U}? The appropriate test for a {U}-shaped relationship},
  journal={Oxford Bulletin of Economics and Statistics}, volume={72}, number={1},
  pages={109--118}, year={2010}}
@article{lubeamish2004,
  author={Lu, Jane W. and Beamish, Paul W.},
  title={International diversification and firm performance: The {S}-curve hypothesis},
  journal={Academy of Management Journal}, volume={47}, number={4}, pages={598--609}, year={2004}}
@article{contractor2012,
  author={Contractor, Farok J.},
  title={Why do multinational firms exist? A theory note about the effect of multinational
  expansion on performance and recent methodological critiques},
  journal={Global Strategy Journal}, volume={2}, number={4}, pages={318--331}, year={2012}}
@article{marano2016,
  author={Marano, Valentina and Arregle, Jean-Luc and Hitt, Michael A. and Spadafora, Ettore
  and Van Essen, Marc},
  title={Home country institutions and the internationalization-performance relationship:
  A meta-analytic review},
  journal={Journal of Management}, volume={42}, number={5}, pages={1075--1110}, year={2016}}
@book{khanna2010,
  author={Khanna, Tarun and Palepu, Krishna G.},
  title={Winning in Emerging Markets: A Road Map for Strategy and Execution},
  publisher={Harvard Business Press}, year={2010}}
@book{north1990,
  author={North, Douglass C.}, title={Institutions, Institutional Change and Economic Performance},
  publisher={Cambridge University Press}, year={1990}}
@article{cameron2015,
  author={Cameron, A. Colin and Miller, Douglas L.},
  title={A practitioner's guide to cluster-robust inference},
  journal={Journal of Human Resources}, volume={50}, number={2}, pages={317--372}, year={2015}}
@article{aterido2011,
  author={Aterido, Reyes and Hallward-Driemeier, Mary and Pag\'es, Carmen},
  title={Big constraints to small firms' growth? Business environment and employment growth across firms},
  journal={Economic Development and Cultural Change}, volume={59}, number={3}, pages={609--647}, year={2011}}
@article{richard2009,
  author={Richard, Pierre J. and Devinney, Timothy M. and Yip, George S. and Johnson, Gerry},
  title={Measuring organizational performance: Towards methodological best practice},
  journal={Journal of Management}, volume={35}, number={3}, pages={718--804}, year={2009}}
\end{filecontents*}

\title{The Forced Internationalization Penalty: A Monotone-Negative
Boundary Condition of the Internationalization--Performance Relationship in
Pacific Small Island Developing States}
\author{Author details withheld for blind review}
\date{Working draft}

\begin{document}
\maketitle

\begin{abstract}
\noindent The inverted-U is the dominant template for the
internationalization--performance (I--P) relationship, but its applicability
to the most structurally constrained economies is untested. We examine seven
Pacific Small Island Developing States (SIDS) using World Bank Enterprise
Survey establishments and find that the canonical inverted-U does not hold;
instead, internationalization and firm performance are related
\emph{monotonically negatively}. In the main specification with country and
year fixed effects ($N = 959$ firms; $13.8\%$ exporters; mean export intensity
$4.8\%$), the linear export-intensity coefficient is
$\hat{\beta}(\mathrm{FSTS}) = -1.339$ ($SE = 0.386$, $p < 0.001$). Adding a
quadratic term yields no significant curvature, confirming a monotone decline
with no turning point in the observed range. In the very small exporters-only
subsample ($N = 26$) the coefficient remains negative
($\hat{\beta} = -1.176$) but is statistically insignificant ($p = 0.130$),
consistent with low power rather than a sign reversal. We name this pattern
the \emph{Forced Internationalization Penalty} (FIP): firms in micro-economies
export not from competitive advantage but because the domestic market cannot
sustain operations, while extreme logistics costs and thin institutional
support erode returns. FIP is, to our knowledge, the first named
monotone-negative boundary condition of the I--P relationship for the Pacific.\\[4pt]
\textbf{Keywords:} internationalization; firm performance; small island
developing states; boundary condition; forced internationalization; Pacific
\end{abstract}

\section{Introduction}
The inverted-U has become a near-default expectation for the I--P
relationship \citep{lubeamish2004,contractor2012}. Yet the universe of firms
underlying that consensus is dominated by larger emerging and advanced
economies; the smallest, most peripheral economies are almost absent from the
evidence base. This omission matters theoretically: if internationalization
is partly \emph{forced} by domestic market insufficiency rather than chosen on
the basis of capability, the benefits-then-costs logic that produces an
inverted-U may not apply at all.

We study seven Pacific SIDS---a setting that maximizes structural constraint:
minute domestic markets, extreme geographic dispersion, and thin institutional
support. The gap we address is the absence of any I--P evidence, and any
theory, for this boundary. Our contribution is to document and name the
\emph{Forced Internationalization Penalty} (FIP): a robust monotone-negative
I--P relationship that inverts the standard prediction, and to specify the
structural mechanism that generates it.

\section{Theory and Hypotheses}
In standard accounts firms internationalize to exploit a transferable
advantage, and the inverted-U arises because expansion benefits eventually
give way to coordination and foreignness costs
\citep{lubeamish2004,marano2016}. Institutional and market-size conditions
bound this logic \citep{north1990,khanna2010}. In SIDS three structural
features invert it: (1) domestic markets too small to sustain operations,
forcing exporting irrespective of advantage; (2) extreme logistics and
market-access costs from remoteness and dispersion; and (3) weak institutional
and export-capability support. Where exporting is a survival necessity rather
than an advantage-driven choice, higher internationalization should map onto
\emph{lower} performance throughout the observed range.

\begin{description}
\item[H1b (FIP).] In Pacific SIDS the I--P relationship is monotonically
negative, with no interior turning point.
\end{description}

\section{Data and Methods}
The data are World Bank Enterprise Surveys (WBES), establishment-level surveys
under a standardized stratified design \citep{aterido2011}, treated as
repeated cross-sections \citep{deaton1985,verbeek2008}. The SIDS analytic
sample comprises seven Pacific economies (ICRV Group VI), $N = 959$ firms, of
which only $N_{\text{exporters}} = 132$ ($13.8\%$) export; mean export
intensity is $0.048$ ($4.8\%$). [These descriptives are REAL.] Labour
productivity is $\ln(\text{sales}/\text{employees})$, standardized within
country--year \citep{richard2009}; export intensity (FSTS) is the export share
of sales. Following the dissertation's pooled-cross-section strategy, models
carry country and year fixed effects, identifying the relationship from
within-country variation; standard errors are robust and clustered by economy
\citep{cameron2015}. We estimate a linear model (M1) and a quadratic model
(M2) and apply the Lind--Mehlum test for any interior extremum
\citep{lindmehlum2010}.

\section{Results}
\subsection{Monotone-negative I--P (H1b)}
Model M1 (country and year fixed effects) gives a significant negative export
coefficient:
\[
\hat{\beta}(\mathrm{FSTS}) = -1.339,\qquad SE = 0.386,\qquad p < 0.001.
\]
Higher internationalization is associated with \emph{lower} firm performance
in SIDS. Adding FSTS$^2$ (model M2) produces no significant curvature,
confirming a monotone decline with no turning point in the data range. In the
exporters-only subsample ($N = 26$, very small) the coefficient stays negative
but is insignificant, consistent with insufficient power:
\[
\hat{\beta}(\mathrm{FSTS}) = -1.176,\qquad p = 0.130\ (\text{NS}).
\]
[All values in this subsection are REAL.] The dissertation's hypothesis
synthesis additionally reports the coefficient ranging to
$-3.351$ ($p<0.001$) under a year-fixed-effects-only specification, confirming
robustness of the negative sign across fixed-effect structures.

% --------------------------------------------------------------------
\begin{table}[H]
\centering
\caption{SIDS I--P specifications (full coefficients).}
\begin{tabularx}{\textwidth}{Xccc}
\toprule
 & M1 (full sample) & M2 (quadratic) & Exporters-only \\
\midrule
\multicolumn{4}{c}{\itshape [Table P8-1 -- full coefficients from the authors'
original output. Populate only from locked estimates.]}\\
FSTS            & $-1.339^{***}$ & --- & $-1.176$ \\
\ (SE)          & $(0.386)$ & --- & --- \\
FSTS$^2$        & --- & ns & --- \\
$p$             & $<0.001$ & --- & $0.130$ \\
$N$             & $959$ & $959$ & $26$ \\
Country FE / Year FE & Yes / Yes & Yes / Yes & Yes / Yes \\
\bottomrule
\end{tabularx}
\end{table}
% --------------------------------------------------------------------

\section{Discussion and Conclusion}
The SIDS evidence identifies a genuine boundary condition of the I--P
literature. Where the inverted-U presumes internationalization is an
advantage-driven choice, SIDS firms internationalize because their domestic
markets are too small to sustain them; under remoteness-driven logistics costs
and thin institutional support, additional internationalization erodes rather
than builds performance. We label this the Forced Internationalization Penalty
(FIP). FIP complements, rather than contradicts, the broader contingent view
of the I--P relationship: it is the limiting case where institutional and
market-size constraints are so severe that the ascending arm of the curve
disappears entirely.

The principal limitations are statistical power and design. The exporter
subsample is tiny ($N = 26$), so exporters-only estimates are
underpowered; and repeated cross-sections preclude dynamic inference. These
caveats notwithstanding, the negative sign is consistent across fixed-effect
structures. Future work should extend FIP to other micro-economy contexts and
test policy levers (logistics subsidies, export-capability programs) that
might restore an ascending arm.

\bibliography{p8refs}

\end{document}
